\documentclass[11pt]{article}
\usepackage[margin=1in]{geometry}
\usepackage{amsmath,amssymb,amsthm}
\usepackage{hyperref}
\usepackage{xurl}
\sloppy
\let\oldus\_
\renewcommand{\_}{\oldus\allowbreak}

\newtheorem{theorem}{Theorem}
\newtheorem{lemma}[theorem]{Lemma}
\theoremstyle{definition}
\newtheorem{definition}[theorem]{Definition}
\theoremstyle{remark}
\newtheorem{remark}[theorem]{Remark}

\title{A disproof of Conjecture 00000000314\\[2pt]
\large Every $M(\sqrt p)$ is at least $2\sqrt 2$, so $\{M(\sqrt p)\}$ is not dense in $[\sqrt5,3)$}
\author{Jackmeson1}
\date{}

\begin{document}
\maketitle

\section{The conjecture}
Conjecture 00000000314 reads (English text, verbatim): ``Definition: The Markov constant
$M(\alpha) = \limsup 1/\|q\alpha\|$ measures the quality of worst approximation. Conjecture: The set
$\{M(\sqrt p) : p \text{ prime}\}$ is dense in the half-open interval from $\sqrt5$ to $3$; and the
density follows from the transition statistics of the continued-fraction partial quotients of square
roots of primes.'' The Chinese text gives the same definition and the same interval.

The statement is a conjunction. We refute its first conjunct (density); the second conjunct
(``the density follows from \dots'') presupposes the first and falls with it.

\section{Definitions and readings}
\begin{definition}
For $x\in\mathbb R$, $\|x\| = \inf_{n\in\mathbb Z}|x-n|$ is the distance from $x$ to the nearest
integer (Lean: \texttt{distNearestInt}). Limits superior are taken over positive integers
$q\to\infty$ and values lie in $[0,\infty]$, with the convention $1/0=\infty$.
\end{definition}
We treat two readings of $M$.
\begin{itemize}
\item \textbf{Standard reading} (Lean: \texttt{markovConst}):
  $M(\alpha) = \limsup_{q\to\infty} \dfrac{1}{q\,\|q\alpha\|}$ $\bigl(= 1/\liminf_q q\|q\alpha\|\bigr)$.
  This is the usual Lagrange value of $\alpha$; its values on irrationals form the Lagrange spectrum,
  whose smallest element is $\sqrt5$, which is why the text's interval starts at $\sqrt5$.
\item \textbf{Literal reading} (Lean: \texttt{markovConstLiteral}):
  $M(\alpha) = \limsup_{q\to\infty} \dfrac{1}{\|q\alpha\|}$, exactly as printed (no factor $q$).
\end{itemize}
``$S$ is dense in $I=[\sqrt5,3)$'' is formalized in two ways, and both are refuted:
(T) $I\subseteq\overline S$ (Lean: \texttt{DenseInTop}); (I) for all $a<b$ in $I$ there is
$s\in S$ with $a<s<b$ (Lean: \texttt{DenseInIntervals}). Both are taken in $[0,\infty]$
(Lean type \texttt{ENNReal}), and $S=\{M(\sqrt p): p\text{ prime}\}$ (Lean: \texttt{valueSet}).

\section{The proof}
\begin{lemma}[Pell estimate; Lean: \texttt{pell\_bound}]
Let $d\ge1$ and let $x,y$ be positive integers with $x^2-dy^2=1$. Then
$\|y\sqrt d\|\cdot 2y\sqrt d\le 1$.
\end{lemma}
\begin{proof}
Put $s=\sqrt d$. Since $(x-ys)(x+ys)=x^2-dy^2=1$ and $x+ys>0$, we get $x-ys>0$, so $x>ys$.
Hence $\|ys\|\le|ys-x| = x-ys$, and
$\|ys\|\cdot 2ys \le (x-ys)\cdot 2ys\le (x-ys)(x+ys)=1$.
\end{proof}

\begin{lemma}[Lean: \texttt{pell\_large}, \texttt{exists\_good\_q}]
For every prime $p$ and every $N$ there is an integer $q> N$ with
$\|q\sqrt p\|\cdot 2q\sqrt p\le1$.
\end{lemma}
\begin{proof}
A prime is not a perfect square, so Pell's equation $x^2-py^2=1$ has a fundamental solution $a$
(Mathlib: \texttt{Pell.IsFundamental.exists\_of\_not\_isSquare}). The map $n\mapsto y(a^n)$ is
strictly increasing on $\mathbb Z$ (Mathlib: \texttt{Pell.IsFundamental.y\_strictMono}) with
$y(a^0)=0$, so $y(a^n)\ge n$ for $n\in\mathbb N$; and $x(a^n)>0$
(Mathlib: \texttt{Pell.Solution\_1.x\_zpow\_pos}; in Mathlib the subscript 1 is a Unicode subscript). Take $q=y(a^{N+1})$ and apply Lemma 1.
\end{proof}

\begin{theorem}[Lean: \texttt{markovConst\_sqrt\_prime\_ge}, \texttt{markovConstLiteral\_sqrt\_prime}]
For every prime $p$: (a) $M(\sqrt p)\ge 2\sqrt p$ in the standard reading;
(b) $M(\sqrt p)=\infty$ in the literal reading.
\end{theorem}
\begin{proof}
(a) For the $q$ of Lemma 2, $q\|q\sqrt p\|\cdot 2\sqrt p\le1$, i.e.\ $1/(q\|q\sqrt p\|)\ge2\sqrt p$.
Such $q$ exist beyond every $N$, so the $\limsup$ is at least $2\sqrt p$.
(b) For every real $B$, apply Lemma 2 with $N\ge B$: then $1/\|q\sqrt p\|\ge 2q\sqrt p\ge q\ge B$,
because $\sqrt p\ge1$. Thus the $\limsup$ is at least every real $B$, so it is $\infty$.
\end{proof}

\begin{theorem}[Lean: \texttt{not\_dense\_of\_ge}, \texttt{conjecture314\_false}]
In both readings of $M$, the set $S=\{M(\sqrt p):p\text{ prime}\}$ is dense in $[\sqrt5,3)$ in
neither sense (T) nor sense (I). In fact $S\cap(\sqrt5,2\sqrt2)=\emptyset$ while
$\sqrt5<2\sqrt2<3$.
\end{theorem}
\begin{proof}
By Theorem 3, every element of $S$ is $\ge 2\sqrt p\ge2\sqrt2$ (standard) or $=\infty$ (literal), so
$S\subseteq[2\sqrt2,\infty]$, a closed set. Also $\sqrt5<2\sqrt2=\sqrt8<\sqrt9=3$.
(T): $\sqrt5\in[\sqrt5,3)$ but $\sqrt5\notin\overline S\subseteq[2\sqrt2,\infty]$.
(I): take $a=\sqrt5$, $b=2\sqrt2$, both in $[\sqrt5,3)$; no $s\in S$ satisfies $s<b$.
\end{proof}

\begin{remark}
For context only (not used in the proof): Wikipedia's article on the Markov spectrum states that
``the initial part of the Markov and Lagrange spectrum lying in the interval $[\sqrt5, 3)$ are both
equal and they are a discrete set''
(\url{https://en.wikipedia.org/wiki/Markov_spectrum}). The argument above is self-contained and
gives the stronger statement that no $M(\sqrt p)$ lies below $2\sqrt2$.
\end{remark}

\section{Lean formalization}
File \texttt{lean/Conjecture314/Basic.lean} (Lean 4, Mathlib \texttt{v4.33.1}, only import
\texttt{Mathlib}). Main theorem:
\begin{verbatim}
theorem conjecture314_false :
    (Not (DenseInTop (valueSet markovConst) targetInterval) /\
      Not (DenseInIntervals (valueSet markovConst) targetInterval)) /\
    (Not (DenseInTop (valueSet markovConstLiteral) targetInterval) /\
      Not (DenseInIntervals (valueSet markovConstLiteral) targetInterval))
\end{verbatim}
(ASCII rendering: the Lean source writes the negations with the Unicode symbol for
\texttt{Not} and the conjunctions with the Unicode wedge.)
Here \texttt{targetInterval = Set.Ico (ENNReal.ofReal (sqrt 5)) 3}. The full chain of lemmas:
\texttt{distNearestInt\_le}, \texttt{pell\_large}, \texttt{pell\_bound},
\texttt{exists\_good\_q}, \texttt{markovConst\_sqrt\_prime\_ge},
\texttt{markovConstLiteral\_sqrt\_prime}, \texttt{sqrt5\_lt\_two\_sqrt2},
\texttt{two\_sqrt2\_lt\_three}, \texttt{not\_dense\_of\_ge}.

\paragraph{Axiom audit.} \texttt{\#print axioms C314.conjecture314\_false} prints
\texttt{[propext, Classical.choice, Quot.sound]}. There is no \texttt{sorry} and no
\texttt{native\_decide}.

\paragraph{Conventions.} $q$ ranges over $\mathbb N$ with $\limsup$ along \texttt{atTop} (the term
$q=0$ is irrelevant to a $\limsup$); values are in $[0,\infty]$ with $1/0=\infty$; $\sqrt{\cdot}$ is
the nonnegative square root. Using $q\in\mathbb Z$, $|q|\to\infty$ gives the same values, since
$\|-x\|=\|x\|$.

\end{document}
